\chapter{Vladimir Drinfeld --- Fields Medal 1990}\label{ch1990vladimirdrinfeld:vladimir-drinfeld-fields-medal-1990-a-comprehensive-report}

\begin{center}\rule{0.5\linewidth}{0.5pt}\end{center}

\textbf{Prepared:} July 2026 \textbf{Subject:} Vladimir Gershonovich
Drinfeld (b. February 14, 1954, Kharkiv, USSR) \textbf{Position:}
Professor, University of Chicago (since 1999) \textbf{Doctoral Advisor:}
Yuri Manin (Moscow State University) \textbf{Fields Medal Citation:}
\emph{``For his contributions to algebraic geometry, number theory, and
mathematical physics, including the theory of quantum groups and the
geometric Langlands program.''}

\begin{center}\rule{0.5\linewidth}{0.5pt}\end{center}

\section{Introduction}\label{ch1990vladimirdrinfeld:introduction}

Vladimir Drinfeld was awarded the Fields Medal in 1990 for a notable
body of work spanning algebraic geometry, number theory, representation
theory, and mathematical physics. His contributions are among the most
original and influential of the late 20th century, introducing entirely
new mathematical disciplines and substantially reshaping existing ones.

In addition to the Fields Medal, Drinfeld has received the Shaw Prize
(2018) and the Wolf Prize (2023), reflecting the enduring impact of his
ideas across multiple generations of mathematicians.

Drinfeld received the medal at the International Congress of
Mathematicians in Kyoto, Japan, together with Vaughan Jones, Shigefumi
Mori, and Edward Witten --- an exceptionally strong cohort that
reflected the substantial breadth of late 20th century mathematics.

Drinfeld's work introduced several substantial concepts: -
\textbf{Drinfeld modules:} A new class of objects in algebraic number
theory that function as the analogue of elliptic curves in the function
field setting. - \textbf{Quantum groups:} A deformation of the universal
enveloping algebra of a Lie algebra, giving rise to a new field at the
intersection of algebra and mathematical physics. - \textbf{The
geometric Langlands program:} A geometric reformulation of the Langlands
conjectures relating local systems on curves to Hecke eigensheaves on
moduli stacks. - \textbf{Yangians:} A particular class of quantum groups
associated to the Yang--Baxter equation, fundamental in integrable
systems. - \textbf{Chtoukas (shtukas):} A generalization of Drinfeld
modules that provides the geometric framework for the Langlands program
over function fields. - \textbf{The Drinfeld double:} A construction
that produces quasi-triangular Hopf algebras from bialgebras, central to
quantum group theory.

The range and depth of Drinfeld's work is notable: he proved the
first non-abelian case of the Langlands conjectures, introduced the
theory of quantum groups that revolutionized representation theory,
co-founded the geometric Langlands program that continues to drive
research in algebraic geometry and quantum field theory, and developed
fundamental tools in integrable systems and representation theory that
bear his name.

\begin{center}\rule{0.5\linewidth}{0.5pt}\end{center}

\section{Biographical Background}\label{ch1990vladimirdrinfeld:biographical-background}

\subsection{Early Life}\label{ch1990vladimirdrinfeld:early-life}

Vladimir Drinfeld was born in 1954 in Kharkiv, Ukraine, then part of the
Soviet Union. He was a child prodigy, winning a gold medal at the
International Mathematical Olympiad in 1969 at the age of 15, achieving
a perfect score. His notable talent was recognized early, and he
was admitted to a special mathematics program.

\subsection{Education}\label{ch1990vladimirdrinfeld:education}

Drinfeld entered Moscow State University in 1969 and graduated in 1974.
He completed his PhD under the supervision of Yuri Manin, one of the
leading mathematicians of the Soviet era, receiving his degree in 1978.
His PhD work already contained the seeds of his substantial
contributions to number theory.

\subsection{Academic Career}\label{ch1990vladimirdrinfeld:academic-career}

\begin{itemize}
\tightlist
\item
 \textbf{1974--1978:} PhD student, Moscow State University, under Yuri
 Manin.
\item
 \textbf{1978--1990:} Researcher, Institute for Low Temperature Physics
 and Engineering, Kharkiv. Despite being isolated from the main
 mathematical centers, Drinfeld produced his most groundbreaking work
 during this period.
\item
 \textbf{1990:} Fields Medal (shared with Vaughan Jones, Shigefumi
 Mori, Edward Witten) at ICM Kyoto.
\item
 \textbf{1990--1998:} Continued research at the Institute for Low
 Temperature Physics and Engineering.
\item
 \textbf{1999--present:} Professor, Department of Mathematics,
 University of Chicago.
\item
 \textbf{2016:} Elected to the National Academy of Sciences (USA).
\item
 \textbf{2018:} Shaw Prize in Mathematical Sciences (shared with
 Alexander Beilinson).
\item
 \textbf{2023:} Wolf Prize in Mathematics.
\end{itemize}

\begin{center}\rule{0.5\linewidth}{0.5pt}\end{center}

\section{Drinfeld Modules}\label{ch1990vladimirdrinfeld:drinfeld-modules}

\subsection{Motivation}\label{ch1990vladimirdrinfeld:motivation}

Drinfeld modules arise from the analogy between number fields and
function fields --- two parallel worlds in arithmetic geometry. Just as
elliptic curves (one-dimensional abelian varieties) are the fundamental
objects used to study number fields via their Galois representations,
Drinfeld modules are used to study function fields over finite fields.

\subsection{Definition}\label{ch1990vladimirdrinfeld:definition}

Let \(F\) be a function field in one variable over a finite field
\(\mathbb{F}_q\). Let \(A\) be the ring of functions regular away from a
fixed place \(\infty\) --- this plays the role of \(\mathbb{Z}\) in the
number field setting. A \textbf{Drinfeld module} of rank \(r\) over an
\(A\)-algebra \(R\) is a homomorphism:

\[
\phi : A \to R\{\tau\}, \quad a \mapsto \phi_a
\]

where \(R\{\tau\}\) is the ring of Ore polynomials over \(R\) with the
commutation relation \(\tau r = r^q \tau\). The homomorphism satisfies a
degree condition:

\[
\deg_\tau \phi_a = r \cdot \deg_{A} a
\]

for all \(a \in A\). The integer \(r \geq 1\) is the rank of the
Drinfeld module.

\subsection{Analogy with Elliptic
Curves}\label{ch1990vladimirdrinfeld:analogy-with-elliptic-curves}

Drinfeld modules play the same role in function field arithmetic that
elliptic curves and abelian varieties play in number field arithmetic: -
The Tate module of a Drinfeld module gives Galois representations. - The
analog of the modular curve is the Drinfeld modular variety. - The
analog of complex multiplication is the theory of Drinfeld modules with
extra endomorphisms.

\subsection{Significance}\label{ch1990vladimirdrinfeld:significance}

Drinfeld modules were used by Drinfeld to achieve several fundamental
results: - Proving the Langlands conjecture for \(GL(2)\) over function
fields --- the first non-abelian case ever proved of any Langlands
conjecture. - Constructing explicit class field theory for function
fields, providing an analytic description of the maximal abelian
extension analogous to the classical theory of complex multiplication. -
Studying the structure of Galois representations arising from function
fields and their local properties. - Laying the foundation for the
theory of chtoukas (shtukas), which later enabled Laurent Lafforgue
(Fields Medal 2002) to prove the Langlands conjecture for \(GL(r)\) over
function fields. - Providing a new perspective on the analogy between
number fields and function fields that continues to inspire research in
arithmetic geometry.

\begin{center}\rule{0.5\linewidth}{0.5pt}\end{center}

\section{The Langlands Program}\label{ch1990vladimirdrinfeld:the-langlands-program}

\subsection{Background}\label{ch1990vladimirdrinfeld:background}

The \textbf{Langlands program}, formulated by Robert Langlands in the
1960s, is one of the most far-reaching web of conjectures in
mathematics. It relates three central objects: Galois representations
(from number theory), automorphic forms (from harmonic analysis), and
algebraic varieties (from algebraic geometry). The central prediction is
that every Galois representation should correspond to an automorphic
representation, known as the Langlands correspondence.

\subsection{Drinfeld's Contribution}\label{ch1990vladimirdrinfeld:drinfelds-contribution}

Drinfeld proved the \textbf{Langlands conjecture for \(GL(2)\) over
function fields} in a series of papers published in the 1970s and early
1980s. This was the first complete proof of a non-abelian case of the
Langlands conjectures --- a landmark achievement that demonstrated the
power of new geometric methods in number theory.

\subsection{The Method}\label{ch1990vladimirdrinfeld:the-method}

Drinfeld's proof was a masterpiece of mathematical ingenuity, combining:
- \textbf{Drinfeld modules} to construct the necessary Galois
representations with prescribed local properties. - \textbf{Modular
varieties} of Drinfeld modules, which serve as the function field
analogues of modular curves and provide the geometric setting for the
correspondence. - \textbf{The trace formula} to match eigenvalues on
both sides of the correspondence, establishing the bijection between
automorphic and Galois representations. - \textbf{Chtoukas} (shtukas),
which Drinfeld introduced as a generalization of Drinfeld modules for
this purpose.

\subsection{Impact}\label{ch1990vladimirdrinfeld:impact}

This work laid the foundation for the \textbf{geometric Langlands
program} and directly influenced: - Lafforgue's proof of the Langlands
conjecture for \(GL(n)\) over function fields (Fields Medal 2002). - The
development of the theory of chtoukas as a central tool in arithmetic
geometry. - Ngô Bảo Châu's proof of the fundamental lemma (Fields Medal
2010), which also arose from the function field Langlands program.

\begin{center}\rule{0.5\linewidth}{0.5pt}\end{center}

\section{The Geometric Langlands
Program}\label{ch1990vladimirdrinfeld:the-geometric-langlands-program}

\subsection{The Correspondence}\label{ch1990vladimirdrinfeld:the-correspondence}

Drinfeld, together with Alexander Beilinson in the 1980s, formulated the
\textbf{geometric Langlands correspondence}, which replaces number
fields with algebraic curves over \(\mathbb{C}\). This geometric
reformulation makes the Langlands conjectures accessible to the powerful
tools of algebraic geometry and representation theory. The
correspondence relates: - \textbf{Local systems} on a smooth projective
curve \(X\) over \(\mathbb{C}\) (the geometric analogue of Galois
representations). - \textbf{Hecke eigensheaves} on the moduli stack
\(\text{Bun}_G\) of principal \(G\)-bundles on \(X\) (the geometric
analogue of automorphic forms).

\subsection{The Statement}\label{ch1990vladimirdrinfeld:the-statement}

Let \(G\) be a reductive algebraic group over \(\mathbb{C}\) with
Langlands dual group \(^L G\), and let \(X\) be a smooth projective
curve. The geometric Langlands conjecture states an equivalence of
derived categories:

\[
\text{Shv}_{G}(\text{Bun}_G) \cong \text{Loc}_{^L G}(X)
\]

where \(\text{Bun}_G\) is the moduli stack of \(G\)-bundles on \(X\) and
\(\text{Loc}_{^L G}(X)\) is the moduli stack of flat \(^L G\)-bundles
(local systems) on \(X\). This equivalence intertwines the Hecke
operators on the left with the restriction functors on the right.

\subsection{Significance}\label{ch1990vladimirdrinfeld:significance-1}

The geometric Langlands program has become one of the most active and
influential areas of contemporary mathematics, connecting: -
\textbf{Representation theory} of affine Lie algebras and loop groups. -
\textbf{Quantum field theory}, particularly S-duality in \(N=4\) super
Yang--Mills theory, where the geometric Langlands correspondence appears
as a statement about electric-magnetic duality. - \textbf{Integrable
systems}, including the Hitchin system and the KdV hierarchy. -
\textbf{Mirror symmetry}, where it relates A-branes and B-branes on dual
moduli spaces. - \textbf{The theory of automorphic sheaves}, which are
the geometric version of automorphic forms.

\begin{center}\rule{0.5\linewidth}{0.5pt}\end{center}

\section{Quantum Groups}\label{ch1990vladimirdrinfeld:quantum-groups}

\subsection{Introduction}\label{ch1990vladimirdrinfeld:introduction-1}

\textbf{Quantum groups} are deformations of the universal enveloping
algebra \(U(\mathfrak{g})\) of a Lie algebra \(\mathfrak{g}\). Despite
their name, they are not groups but rather Hopf algebras --- algebraic
structures that generalize the algebra of functions on a group. Drinfeld
introduced quantum groups in his 1986 ICM address, instantly creating a
new field of mathematics.

\subsection{Definition}\label{ch1990vladimirdrinfeld:definition-1}

The quantum group \(U_q(\mathfrak{g})\) associated to a semi-simple Lie
algebra \(\mathfrak{g}\) is a Hopf algebra over \(\mathbb{C}(q)\)
generated by elements \(E_i, F_i, K_i^{\pm 1}\) for
\(i = 1, \ldots, \text{rank}(\mathfrak{g})\), satisfying:

\[
K_i K_j = K_j K_i, \quad K_i K_i^{-1} = 1
\]

\[
K_i E_j K_i^{-1} = q^{a_{ij}} E_j, \quad K_i F_j K_i^{-1} = q^{-a_{ij}} F_j
\]

\[
[E_i, F_j] = \delta_{ij} \frac{K_i - K_i^{-1}}{q - q^{-1}}
\]

together with the quantum Serre relations. As \(q \to 1\), the algebra
\(U_q(\mathfrak{g})\) specializes to the classical universal enveloping
algebra \(U(\mathfrak{g})\).

\subsection{The Hopf Algebra
Structure}\label{ch1990vladimirdrinfeld:the-hopf-algebra-structure}

The key feature that distinguishes quantum groups from ordinary algebras
is their Hopf algebra structure, which includes: - A
\textbf{comultiplication}
\(\Delta: U_q(\mathfrak{g}) \to U_q(\mathfrak{g}) \otimes U_q(\mathfrak{g})\)
encoding how representations tensor. - A \textbf{counit}
\(\varepsilon: U_q(\mathfrak{g}) \to \mathbb{C}\). - An
\textbf{antipode} \(S: U_q(\mathfrak{g}) \to U_q(\mathfrak{g})\), the
analog of the inverse map.

\subsection{The Yang--Baxter
Equation}\label{ch1990vladimirdrinfeld:the-yangbaxter-equation}

Quantum groups were discovered through Drinfeld's study of the
\textbf{Yang--Baxter equation}, a fundamental equation in statistical
mechanics and integrable systems:

\[
R_{12} R_{13} R_{23} = R_{23} R_{13} R_{12}
\]

Here \(R \in U_q(\mathfrak{g}) \otimes U_q(\mathfrak{g})\) is the
universal \(R\)-matrix, which satisfies this equation. Drinfeld showed
that quantum groups equipped with a universal \(R\)-matrix
(quasi-triangular Hopf algebras) provide a systematic way to construct
solutions of the Yang--Baxter equation, unifying many previously ad-hoc
constructions.

\subsection{The Drinfeld Double}\label{ch1990vladimirdrinfeld:the-drinfeld-double}

Drinfeld introduced the \textbf{Drinfeld double} construction, which
produces a quasi-triangular Hopf algebra \(D(A)\) from a bialgebra \(A\)
and its dual \(A^*\). This construction is a fundamental tool in the
theory of quantum groups and provides a universal method for
constructing solutions to the Yang--Baxter equation.

\subsection{Applications}\label{ch1990vladimirdrinfeld:applications}

Quantum groups have revolutionized several areas of mathematics and
theoretical physics: - \textbf{Statistical mechanics:} Exactly solvable
lattice models, including the XXZ and XYZ spin chains. - \textbf{Knot
theory:} Quantum invariants of knots such as the Jones polynomial and
the HOMFLY polynomial, realized as quantum traces of the \(R\)-matrix. -
\textbf{Representation theory:} Crystal bases and canonical bases
(Lusztig's canonical basis), which provide a combinatorial description
of quantum group representations. - \textbf{Conformal field theory:} The
representation theory of affine Lie algebras and the
Knizhnik--Zamolodchikov equations. - \textbf{Noncommutative geometry:}
Quantum groups as examples of noncommutative spaces.

\begin{center}\rule{0.5\linewidth}{0.5pt}\end{center}

\section{Yangians}\label{ch1990vladimirdrinfeld:yangians}

\subsection{Definition}\label{ch1990vladimirdrinfeld:definition-2}

\textbf{Yangians} are a class of quantum groups introduced by Drinfeld
in the 1980s. The Yangian \(Y(\mathfrak{g})\) of a Lie algebra
\(\mathfrak{g}\) is a Hopf algebra deformation of the universal
enveloping algebra \(U(\mathfrak{g}[t])\) of the polynomial current
algebra. They are named after C. N. Yang, who discovered the
Yang--Baxter equation.

\subsection{Structure}\label{ch1990vladimirdrinfeld:structure}

The Yangian \(Y(\mathfrak{g})\) is an infinite-dimensional Hopf algebra
generated by elements \(\{J^{(n)}_i\}\) for \(n \geq 0\) and
\(i = 1, \ldots, \text{rank}(\mathfrak{g})\), satisfying relations that
reflect the classical \(r\)-matrix. The defining relations have the
form:

\[
[J^{(m)}_i, J^{(n)}_j] = f_{ij}^k J^{(m+n)}_k + \text{(quantum corrections)}
\]

where \(f_{ij}^k\) are the structure constants of \(\mathfrak{g}\). The
Yangian has a natural grading by non-negative integers given by the
``level'' \(n\).

\subsection{The Evaluation Map}\label{ch1990vladimirdrinfeld:the-evaluation-map}

Yangians have a crucial property: there is an \textbf{evaluation
homomorphism}:

\[
\text{ev}: Y(\mathfrak{g}) \to U(\mathfrak{g})
\]

which specializes Yangian representations to ordinary Lie algebra
representations. This allows Yangian representations to be understood as
``families'' of Lie algebra representations.

\subsection{Significance}\label{ch1990vladimirdrinfeld:significance-2}

Yangians appear in several fundamental contexts: - \textbf{Integrable
systems:} The XXX spin chain has Yangian symmetry, with the Yangian
governing the structure of conserved quantities. -
\textbf{Representation theory:} The finite-dimensional irreducible
representations of Yangians have been classified and can be described
using Drinfeld polynomials. - \textbf{Geometric Langlands program:}
Yangians appear as the symmetries of the Hitchin system and the
geometric Langlands correspondence for \(GL(n)\). - \textbf{Quantum
inverse scattering method:} Yangians provide the algebraic framework for
the quantum inverse scattering method.

\begin{center}\rule{0.5\linewidth}{0.5pt}\end{center}

\section{The Drinfeld--Sokolov
Reduction}\label{ch1990vladimirdrinfeld:the-drinfeldsokolov-reduction}

\subsection{Introduction}\label{ch1990vladimirdrinfeld:introduction-2}

The \textbf{Drinfeld--Sokolov reduction} (developed jointly with
Vladimir Sokolov) is a powerful method to construct integrable
Hamiltonian systems from affine Lie algebras. It provides a systematic
framework for understanding the KdV hierarchy and its generalizations.

\subsection{The Construction}\label{ch1990vladimirdrinfeld:the-construction}

Given an affine Lie algebra \(\hat{\mathfrak{g}}\) with central
extension, the Drinfeld--Sokolov reduction proceeds in several steps: 1.
One considers the space of currents (connections) on the circle with
values in \(\hat{\mathfrak{g}}\). 2. A set of constraints is imposed on
the currents, reducing the phase space. 3. A gauge fixing procedure
eliminates the redundant degrees of freedom. 4. The resulting reduced
phase space has a Poisson structure isomorphic to the W-algebra
associated to \(\mathfrak{g}\).

\subsection{The KdV Hierarchy}\label{ch1990vladimirdrinfeld:the-kdv-hierarchy}

For \(\mathfrak{g} = \mathfrak{sl}(2)\), the Drinfeld--Sokolov reduction
produces the classic Korteweg--de Vries (KdV) equation:

\[
u_t = 6uu_x - u_{xxx}
\]

For higher-rank Lie algebras, one obtains the generalized KdV
(Gelfand--Dickey) hierarchies.

\subsection{Applications}\label{ch1990vladimirdrinfeld:applications-1}

The Drinfeld--Sokolov reduction has found applications in: - \textbf{KdV
hierarchy and its generalizations:} A unified framework for classical
integrable systems. - \textbf{W-algebras:} Extended conformal algebras
that appear in two-dimensional conformal field theory and the
representation theory of affine Lie algebras. - \textbf{Integrable
systems:} The reduction provides a systematic classification of
integrable Hamiltonian systems associated to Lie algebras. -
\textbf{Quantum field theory:} W-algebras play a central role in the
study of WZW models and their deformations.

\begin{center}\rule{0.5\linewidth}{0.5pt}\end{center}

\section{The Hitchin--Drinfeld
System}\label{ch1990vladimirdrinfeld:the-hitchindrinfeld-system}

\subsection{Definition}\label{ch1990vladimirdrinfeld:definition-3}

The \textbf{Hitchin--Drinfeld system} (now known as the Hitchin
integrable system) is an algebraically completely integrable Hamiltonian
system on the moduli space of Higgs bundles on a Riemann surface. It was
discovered by Nigel Hitchin and independently by Drinfeld. The system is
defined by the Hitchin map:

\[
h: \mathcal{M}_{\text{Higgs}}(G) \to \bigoplus_{i=1}^{\text{rank}(G)} H^0(X, K^{\otimes d_i})
\]

which sends a Higgs bundle \((E, \phi)\) to the coefficients of the
characteristic polynomial of \(\phi\).

\subsection{Relationship to the Langlands
Program}\label{ch1990vladimirdrinfeld:relationship-to-the-langlands-program}

The Hitchin system plays a central role in the geometric Langlands
program: - The Hitchin fibration provides a geometric realization of the
Langlands duality between \(G\) and its Langlands dual \(^L G\). - The
fibers of the Hitchin map are abelian varieties (Jacobians of spectral
curves), giving an abelianized version of the correspondence. - The
Hitchin system is related to the moduli space of flat connections via
non-abelian Hodge theory (the Simpson correspondence). - Mirror symmetry
for Higgs bundles relates the Hitchin system for \(G\) to the Hitchin
system for \(^L G\).

\begin{center}\rule{0.5\linewidth}{0.5pt}\end{center}

\section{The Braid Group and Hecke
Algebras}\label{ch1990vladimirdrinfeld:the-braid-group-and-hecke-algebras}

\subsection{Drinfeld's Work}\label{ch1990vladimirdrinfeld:drinfelds-work}

Drinfeld studied the representation theory of braid groups and Hecke
algebras, using his theory of quantum groups to construct new and
powerful representations. The braid group \(B_n\) on \(n\) strands is
generated by elementary braids \(\sigma_1, \ldots, \sigma_{n-1}\) with
relations:

\[
\sigma_i \sigma_{i+1} \sigma_i = \sigma_{i+1} \sigma_i \sigma_{i+1}, \quad \sigma_i \sigma_j = \sigma_j \sigma_i \text{ for } |i - j| > 1
\]

\subsection{The KZ Connection}\label{ch1990vladimirdrinfeld:the-kz-connection}

The \textbf{Knizhnik--Zamolodchikov (KZ) connection} is a flat
connection on the vector bundle of conformal blocks of the
Wess--Zumino--Witten (WZW) model. It is defined as:

\[
\nabla_i = \frac{\partial}{\partial z_i} - \kappa \sum_{j \neq i} \frac{\Omega_{ij}}{z_i - z_j}
\]

where \(\Omega_{ij}\) is the Casimir operator acting on the \(i\)-th and
\(j\)-th tensor factors. Drinfeld showed a fundamental result: the
monodromy of the KZ connection is given by the quantum \(R\)-matrix of
the corresponding quantum group. This deep connection between conformal
field theory and quantum groups has had far-reaching implications for
both subjects.

\subsection{The Kohno--Drinfeld
Theorem}\label{ch1990vladimirdrinfeld:the-kohnodrinfeld-theorem}

The \textbf{Kohno--Drinfeld theorem} (proved independently by Drinfeld
and Toshitake Kohno) states that the monodromy representation of the
braid group obtained from the KZ connection is equivalent to the
representation obtained from the quantum group \(U_q(\mathfrak{g})\)
when \(q = \exp(\pi i / \kappa)\). This provides a precise
correspondence between conformal field theory and quantum groups.

\begin{center}\rule{0.5\linewidth}{0.5pt}\end{center}

\section{The Work on the Deligne--Lusztig
Theory}\label{ch1990vladimirdrinfeld:the-work-on-the-delignelusztig-theory}

\subsection{Background}\label{ch1990vladimirdrinfeld:background-1}

\textbf{Deligne--Lusztig theory} (developed by Pierre Deligne and George
Lusztig in the 1970s) constructs representations of finite groups of Lie
type, such as \(GL(n, \mathbb{F}_q)\), using the \(l\)-adic cohomology
of algebraic varieties. This theory unified the representation theory of
finite simple groups and led to the classification of their irreducible
representations.

\subsection{Drinfeld's
Contributions}\label{ch1990vladimirdrinfeld:drinfelds-contributions}

Drinfeld contributed several important variants and applications of the
Deligne--Lusztig construction: - Using Drinfeld modules to construct
representations of the general linear group over finite fields,
providing a more explicit geometric realization. - Developing the
connection between Deligne--Lusztig theory and the Langlands program for
function fields. - Extending the theory to higher-rank groups using
chtoukas and their moduli spaces. - Relating the cohomology of
Deligne--Lusztig varieties to the representation theory of quantum
groups at roots of unity.

\begin{center}\rule{0.5\linewidth}{0.5pt}\end{center}

\section{The Work on Formal Groups and Local
Fields}\label{ch1990vladimirdrinfeld:the-work-on-formal-groups-and-local-fields}

\subsection{Formal Groups}\label{ch1990vladimirdrinfeld:formal-groups}

A \textbf{formal group} (or more precisely, a one-parameter formal group
law) is a formal power series \(F(X, Y) \in R[[X, Y]]\) satisfying: -
\(F(X, 0) = X = F(0, X)\) (identity). -
\(F(F(X, Y), Z) = F(X, F(Y, Z))\) (associativity). -
\(F(X, Y) = X + Y + \text{(higher order terms)}\).

Formal groups arise naturally in algebraic number theory through the
theory of local fields, where they encode the structure of the
multiplicative group and the addition law on elliptic curves.

\subsection{Drinfeld's
Contributions}\label{ch1990vladimirdrinfeld:drinfelds-contributions-1}

Drinfeld made important contributions to the theory of formal groups: -
The classification of formal groups over local fields of positive
characteristic, using Drinfeld modules as a tool. - The relationship
between formal groups and Drinfeld modules, showing that Drinfeld
modules can be viewed as a global version of formal groups. - The
application of formal groups to the local Langlands program for function
fields. - The construction of Lubin--Tate type theories using Drinfeld
modules, providing explicit class field theory for function fields.

\begin{center}\rule{0.5\linewidth}{0.5pt}\end{center}

\section{The Work on the Ring of Differential
Operators}\label{ch1990vladimirdrinfeld:the-work-on-the-ring-of-differential-operators}

\subsection{The Ring of Differential
Operators}\label{ch1990vladimirdrinfeld:the-ring-of-differential-operators}

The ring of differential operators \(D(X)\) on a smooth algebraic
variety \(X\) is a non-commutative algebra generated by the ring of
regular functions \(O(X)\) and the vector fields on \(X\). It plays a
central role in representation theory through the Beilinson--Bernstein
localization theorem, which relates representations of Lie algebras to
D-modules on flag varieties.

\subsection{Drinfeld's
Contributions}\label{ch1990vladimirdrinfeld:drinfelds-contributions-2}

Drinfeld studied several aspects of the ring of differential operators:
- The structure of the ring of differential operators on flag varieties
\(G/B\), particularly its center and its relationship to the universal
enveloping algebra. - The relationship between differential operators
and representations of Lie algebras, including the Beilinson--Bernstein
correspondence. - The application of D-modules to the geometric
Langlands program, where D-modules on \(\text{Bun}_G\) replace
automorphic forms. - The connection between differential operators and
quantum groups through the theory of quantum D-modules.

\begin{center}\rule{0.5\linewidth}{0.5pt}\end{center}

\section{The Work on the Trace
Formula}\label{ch1990vladimirdrinfeld:the-work-on-the-trace-formula}

\subsection{The Arthur--Selberg Trace
Formula}\label{ch1990vladimirdrinfeld:the-arthurselberg-trace-formula}

The \textbf{Arthur--Selberg trace formula} is one of the effective
tools in automorphic representation theory. For a reductive group \(G\)
over a global field \(F\), it relates: - \textbf{The spectral side:}
Traces of representations appearing in the space of automorphic forms. -
\textbf{The geometric side:} Orbital integrals over conjugacy classes of
\(G(F)\).

\subsection{Drinfeld's
Modifications}\label{ch1990vladimirdrinfeld:drinfelds-modifications}

Drinfeld developed variants of the trace formula adapted to the function
field setting. His key contributions include: - Using Drinfeld modules
to simplify the geometric side of the trace formula, making the orbital
integrals more computable. - Developing a trace formula for the moduli
spaces of Drinfeld modules, which can be seen as a geometric version of
the Arthur--Selberg trace formula. - Connecting the trace formula to the
cohomology of moduli spaces of chtoukas, which was later used by
Lafforgue in his proof of the Langlands conjectures for \(GL(n)\).

\begin{center}\rule{0.5\linewidth}{0.5pt}\end{center}

\section{The Work on Chtoukas}\label{ch1990vladimirdrinfeld:the-work-on-chtoukas}

\subsection{Definition}\label{ch1990vladimirdrinfeld:definition-4}

\textbf{Chtoukas} (also called shtukas, from the Russian word for
``sheaf'') are a generalization of Drinfeld modules introduced by
Drinfeld in his work on the Langlands program. A chtouka over a curve
\(X\) over a finite field \(\mathbb{F}_q\) is a vector bundle \(E\) on
\(X\) equipped with a Frobenius structure: an isomorphism between \(E\)
and its pullback by the Frobenius map, with possible modifications at a
finite set of points.

\subsection{The Moduli Space}\label{ch1990vladimirdrinfeld:the-moduli-space}

Drinfeld constructed moduli spaces of chtoukas, which play the role of
Shimura varieties in the function field setting. These moduli spaces are
smooth algebraic stacks that support the action of Hecke operators and
admit a natural compactification.

\subsection{Significance}\label{ch1990vladimirdrinfeld:significance-3}

Chtoukas are fundamental tools in modern arithmetic geometry: -
\textbf{Langlands program:} The moduli spaces of chtoukas are used to
construct the geometric realizations of automorphic representations,
enabling the proof of the Langlands conjecture for \(GL(r)\) over
function fields by Lafforgue. - \textbf{Cohomology:} The cohomology of
moduli spaces of chtoukas provides explicit realizations of Galois
representations. - \textbf{Shimura varieties:} Chtoukas can be seen as
the function field analogues of Shimura varieties, connecting the
arithmetic geometry of function fields to the theory of automorphic
forms.

\begin{center}\rule{0.5\linewidth}{0.5pt}\end{center}

\section{The Work on the Elliptic Lie
Algebras}\label{ch1990vladimirdrinfeld:the-work-on-the-elliptic-lie-algebras}

\subsection{Definition}\label{ch1990vladimirdrinfeld:definition-5}

\textbf{Elliptic Lie algebras} are a class of infinite-dimensional
Kac--Moody algebras related to elliptic curves. They are central
extensions (and sometimes twisted versions) of the Lie algebra of maps
from an elliptic curve to a finite-dimensional simple Lie algebra. They
appear in the theory of elliptic integrable systems and in the
representation theory of quantum affine algebras.

\subsection{Drinfeld's
Contributions}\label{ch1990vladimirdrinfeld:drinfelds-contributions-3}

Drinfeld contributed to the theory of elliptic Lie algebras in several
ways: - The representation theory of elliptic Lie algebras, including
the construction of Verma modules and highest-weight representations. -
The construction of integrable systems associated to elliptic curves,
including the elliptic Calogero--Moser system. - The relationship
between elliptic Lie algebras and the theory of quantum groups,
particularly the elliptic quantum groups that arise as deformations of
quantum affine algebras.

\begin{center}\rule{0.5\linewidth}{0.5pt}\end{center}

\section{The Work on the Virasoro
Algebra}\label{ch1990vladimirdrinfeld:the-work-on-the-virasoro-algebra}

\subsection{Definition}\label{ch1990vladimirdrinfeld:definition-6}

The \textbf{Virasoro algebra} is an infinite-dimensional Lie algebra
that plays a central role in conformal field theory and string theory.
It is generated by elements \(\{L_n\}_{n\in\mathbb{Z}}\) and a central
element \(c\) satisfying the commutation relations:

\[
[L_m, L_n] = (m-n)L_{m+n} + \frac{c}{12}(m^3 - m)\delta_{m+n,0}
\]

The Virasoro algebra is the unique central extension of the Lie algebra
of polynomial vector fields on the circle.

\subsection{Drinfeld's
Contributions}\label{ch1990vladimirdrinfeld:drinfelds-contributions-4}

Drinfeld contributed to the theory of the Virasoro algebra: - The
representation theory of the Virasoro algebra, including the
classification of unitary highest-weight representations. - The
connection between Virasoro representations and the KdV hierarchy
through the Drinfeld--Sokolov reduction, where the KdV hierarchy appears
as a reduction of the Virasoro algebra. - The relationship between
Virasoro algebras and conformal field theory, particularly the role of
Virasoro representations in the WZW model and minimal models.

\begin{center}\rule{0.5\linewidth}{0.5pt}\end{center}

\section{The Work on the McKay
Correspondence}\label{ch1990vladimirdrinfeld:the-work-on-the-mckay-correspondence}

\subsection{The Correspondence}\label{ch1990vladimirdrinfeld:the-correspondence-1}

The \textbf{McKay correspondence}, discovered by John McKay in the
1970s, is a notable relationship between finite subgroups of
\(SL(2,\mathbb{C})\) and simply-laced Dynkin diagrams. For each finite
subgroup \(\Gamma \subset SL(2,\mathbb{C})\), the McKay quiver (whose
vertices are irreducible representations of \(\Gamma\)) is the extended
Dynkin diagram of a simply-laced Lie algebra.

\subsection{Drinfeld's
Contributions}\label{ch1990vladimirdrinfeld:drinfelds-contributions-5}

Drinfeld contributed to the McKay correspondence in several ways: -
Relating it to the representation theory of affine Lie algebras through
the theory of Kleinian singularities and their resolutions. - Connecting
the McKay correspondence to the theory of singularities, particularly
the ADE classification of simple singularities. - Extending the
correspondence to finite subgroups of \(SL(n,\mathbb{C})\) and
higher-dimensional singularities. - Relating the McKay correspondence to
the theory of quiver varieties and the geometric Langlands program.

\begin{center}\rule{0.5\linewidth}{0.5pt}\end{center}

\section{The Work on the Riemann--Hilbert
Correspondence}\label{ch1990vladimirdrinfeld:the-work-on-the-riemannhilbert-correspondence}

\subsection{The Correspondence}\label{ch1990vladimirdrinfeld:the-correspondence-2}

The \textbf{Riemann--Hilbert correspondence} is a fundamental result in
algebraic analysis that establishes an equivalence between regular
holonomic D-modules and perverse sheaves on a complex manifold. It
provides a bridge between analysis (differential equations) and topology
(sheaf theory).

\subsection{Drinfeld's
Contributions}\label{ch1990vladimirdrinfeld:drinfelds-contributions-6}

Drinfeld contributed to several aspects of the Riemann--Hilbert
correspondence: - \textbf{Non-abelian Hodge theory for algebraic
curves:} Drinfeld worked on the correspondence between Higgs bundles,
flat connections, and local systems on algebraic curves, a non-abelian
generalization of the classical Hodge theory. - \textbf{Geometric
Langlands:} The Riemann--Hilbert correspondence is a crucial ingredient
in the geometric Langlands program, relating D-modules on
\(\text{Bun}_G\) (the automorphic side) to local systems (the Galois
side). - \textbf{De Rham moduli space:} The moduli space of flat
connections (the de Rham moduli space) and its relationship to the
Dolbeault moduli space of Higgs bundles was studied by Drinfeld in
connection with the geometric Langlands program.

\begin{center}\rule{0.5\linewidth}{0.5pt}\end{center}

\section{The Work on the Tamagawa
Numbers}\label{ch1990vladimirdrinfeld:the-work-on-the-tamagawa-numbers}

\subsection{Definition}\label{ch1990vladimirdrinfeld:definition-7}

\textbf{Tamagawa numbers} are volume measures associated to adelic
groups, introduced by Tsuneo Tamagawa in the 1960s. For a linear
algebraic group \(G\) defined over a global field \(F\), the Tamagawa
number \(\tau(G)\) is defined as the volume of the quotient
\(G(\mathbb{A})/G(F)\) with respect to the Tamagawa measure --- a
canonical invariant measure on the adelic group. Tamagawa numbers play a
central role in the Eichler--Shimura theory and the
Birch--Swinnerton-Dyer conjecture.

\subsection{Drinfeld's
Contributions}\label{ch1990vladimirdrinfeld:drinfelds-contributions-7}

Drinfeld contributed to the theory of Tamagawa numbers in the function
field setting: - Computing Tamagawa numbers for certain linear algebraic
groups over function fields using the theory of Drinfeld modules. -
Relating these computations to the arithmetic of function fields and the
structure of Galois cohomology. - These computations provided crucial
evidence for the general conjectures about Tamagawa numbers of linear
algebraic groups, which later became theorems for function fields.

\begin{center}\rule{0.5\linewidth}{0.5pt}\end{center}

\section{The Work on the Poisson
Geometry}\label{ch1990vladimirdrinfeld:the-work-on-the-poisson-geometry}

\subsection{Poisson Structures}\label{ch1990vladimirdrinfeld:poisson-structures}

A \textbf{Poisson structure} on a smooth manifold \(M\) is a Lie bracket
\(\{\cdot, \cdot\}\) on the algebra of smooth functions \(C^\infty(M)\)
satisfying the Leibniz rule:

\[
\{f, gh\} = \{f, g\}h + g\{f, h\}
\]

Poisson geometry is the geometric study of classical mechanics: the
phase space of a classical mechanical system is a Poisson manifold, with
the Poisson bracket encoding the dynamics.

\subsection{Drinfeld's
Contributions}\label{ch1990vladimirdrinfeld:drinfelds-contributions-8}

Drinfeld contributed to Poisson geometry in several fundamental ways: -
The \textbf{Drinfeld--Sokolov reduction} as a Poisson reduction, showing
that the KdV hierarchy can be understood as a reduced Hamiltonian system
on an affine Lie algebra. - \textbf{Poisson--Lie groups} and their
classification. A Poisson--Lie group is a Lie group equipped with a
Poisson structure compatible with the group multiplication. Drinfeld
classified all Poisson--Lie structures on simple Lie groups. - The
\textbf{quantization of Poisson manifolds} via quantum groups, where
\(U_q(\mathfrak{g})\) can be seen as a quantization of the Poisson--Lie
structure on \(G\). - The theory of Poisson--Lie duality and its
relationship to the Langlands correspondence.

\begin{center}\rule{0.5\linewidth}{0.5pt}\end{center}

\section{The Work on the Fusion
Categories}\label{ch1990vladimirdrinfeld:the-work-on-the-fusion-categories}

\subsection{Definition}\label{ch1990vladimirdrinfeld:definition-8}

\textbf{Fusion categories} are semisimple tensor categories with
finitely many simple objects and a finite-dimensional space of
morphisms. They arise naturally in the representation theory of quantum
groups and finite groups, as well as in conformal field theory and
topological quantum field theory. The key data of a fusion category
includes: - The set of simple objects and their fusion rules. - The
associativity constraints (F-matrices). - The braiding (for braided
fusion categories).

\subsection{Drinfeld's
Contributions}\label{ch1990vladimirdrinfeld:drinfelds-contributions-9}

Drinfeld contributed to the theory of fusion categories in several
important ways: - \textbf{Classification:} The classification of fusion
categories associated to quantum groups \(U_q(\mathfrak{g})\) at roots
of unity, where the representation theory becomes semisimple with
finitely many irreducible representations. - \textbf{Modular tensor
categories:} The relationship between fusion categories and modular
tensor categories, which describe the representation theory of loop
groups and conformal field theories. - \textbf{Topological quantum field
theory:} The connection between fusion categories and TQFT, where
modular tensor categories give rise to 3-dimensional TQFTs via the
Reshetikhin--Turaev construction. - \textbf{The Drinfeld center:} The
center of a fusion category is a braided fusion category that encodes
the algebraic structure of the category. The Drinfeld center
construction provides a way to produce modular tensor categories from
spherical fusion categories. - \textbf{Applications:} Fusion categories
have applications in quantum computation, particularly in the theory of
topological quantum computers where anyons are described by braided
fusion categories.

\begin{center}\rule{0.5\linewidth}{0.5pt}\end{center}

\section{The Work on the Quantum
Duality}\label{ch1990vladimirdrinfeld:the-work-on-the-quantum-duality}

\subsection{Drinfeld's Theorem}\label{ch1990vladimirdrinfeld:drinfelds-theorem}

Drinfeld proved the \textbf{quantum duality theorem}, a fundamental
result that establishes a precise relationship between the two main
approaches to quantum groups: - The deformation \(U_q(\mathfrak{g})\) of
the universal enveloping algebra, which is a noncommutative,
non-cocommutative Hopf algebra. - The quantized function algebra
\(\mathbb{C}_q[G]\), which is a noncommutative but cocommutative Hopf
algebra.

The theorem states that \(U_q(\mathfrak{g})\) and \(\mathbb{C}_q[G]\)
are dual Hopf algebras in a precise sense: each is the Hopf algebra dual
of the other. This duality is the quantum analogue of the duality
between the universal enveloping algebra \(U(\mathfrak{g})\) and the
algebra of regular functions \(\mathbb{C}[G]\) on the corresponding
algebraic group \(G\).

\subsection{Consequences}\label{ch1990vladimirdrinfeld:consequences}

The quantum duality theorem has important consequences: - It shows that
the two approaches to quantum group theory --- the deformation of
universal enveloping algebras and the quantization of function algebras
--- are equivalent at the level of Hopf algebras. - It provides a
framework for transferring constructions between \(U_q(\mathfrak{g})\)
and \(\mathbb{C}_q[G]\), such as representations and cohomology. - It is
essential for the representation theory of quantum groups at roots of
unity, where the duality becomes especially rich. - It leads to the
theory of quantum groups as quantum symmetry groups of noncommutative
spaces.

\subsection{Applications}\label{ch1990vladimirdrinfeld:applications-2}

The quantum duality has been applied to: - The classification of quantum
group representations at roots of unity, where the duality provides a
correspondence between representations of \(U_q(\mathfrak{g})\) and
\(\mathbb{C}_q[G]\). - The construction of quantum invariants of
3-manifolds via the Reshetikhin--Turaev construction. - The theory of
quantum homogeneous spaces and quantum principal bundles. - The study of
quantum subgroups and quantum symmetric pairs.

\begin{center}\rule{0.5\linewidth}{0.5pt}\end{center}

\section{The Work on the Gelfand--Tsetlin
Basis}\label{ch1990vladimirdrinfeld:the-work-on-the-gelfandtsetlin-basis}

\subsection{Definition}\label{ch1990vladimirdrinfeld:definition-9}

The \textbf{Gelfand--Tsetlin basis} (also called the Gelfand--Tsetlin
patterns) is a canonical basis for the finite-dimensional irreducible
representations of \(GL(n)\) indexed by patterns of integers. A
Gelfand--Tsetlin pattern is a triangular array:

\[
\begin{pmatrix}
m_{1,n} & m_{2,n} & \cdots & m_{n,n} \\
& m_{1,n-1} & \cdots & m_{n-1,n-1} \\
& & \ddots & \\
& & & m_{1,1}
\end{pmatrix}
\]

satisfying interlacing conditions between successive rows.

\subsection{Drinfeld's
Contributions}\label{ch1990vladimirdrinfeld:drinfelds-contributions-10}

Drinfeld studied several aspects of Gelfand--Tsetlin theory: -
\textbf{Quantum deformation:} The quantum Gelfand--Tsetlin basis for
\(U_q(GL(n))\), providing an explicit combinatorial description of
quantum group representations. - \textbf{Yangians:} The relationship
between Gelfand--Tsetlin bases and Yangians, showing that the basis can
be understood in terms of the representation theory of Yangians. -
\textbf{Canonical bases:} The connection to Lusztig's canonical bases
and Kashiwara's crystal bases, which provide a deeper combinatorial
understanding of quantum group representations. - \textbf{Geometric
realization:} The Gelfand--Tsetlin basis can be realized geometrically
via the theory of quiver varieties and the geometric Langlands program.

\begin{center}\rule{0.5\linewidth}{0.5pt}\end{center}

\section{Influence and Legacy}\label{ch1990vladimirdrinfeld:influence-and-legacy}

\subsection{Impact on Number Theory}\label{ch1990vladimirdrinfeld:impact-on-number-theory}

Drinfeld's work on the Langlands program for function fields opened an
entirely new direction in number theory. His use of Drinfeld modules and
chtoukas provided the geometric tools that enabled Lafforgue to prove
the Langlands conjecture for \(GL(n)\) over function fields (Fields
Medal 2002) and influenced Ngô Bảo Châu's proof of the fundamental lemma
(Fields Medal 2010).

\subsection{Impact on Representation
Theory}\label{ch1990vladimirdrinfeld:impact-on-representation-theory}

The theory of quantum groups, which Drinfeld created, revolutionized
representation theory. Quantum groups and their representations are now
a standard tool in: - The representation theory of Lie algebras and
affine Lie algebras. - The theory of knot and link invariants. - The
theory of tensor categories and fusion categories. - The study of
integrable systems.

\subsection{Impact on Mathematical
Physics}\label{ch1990vladimirdrinfeld:impact-on-mathematical-physics}

Drinfeld's contributions to mathematical physics are equally significant: -
The Yang--Baxter equation and quantum groups provide the algebraic
framework for exactly solvable models in statistical mechanics. - The
geometric Langlands program has been connected to S-duality in \(N=4\)
super Yang--Mills theory by Kapustin and Witten, revealing deep
connections between mathematics and quantum field theory. - The
Drinfeld--Sokolov reduction unifies the theory of integrable PDEs.

\section{Awards and Recognition}\label{ch1990vladimirdrinfeld:awards-and-recognition}

\textbf{1990:} - Fields Medal (International Mathematical Union) ---
awarded at the ICM in Kyoto, Japan, for his contributions to algebraic
geometry, number theory, and mathematical physics.

\textbf{1992:} - Corresponding Member, Ukrainian Academy of Sciences.

\textbf{2016:} - Elected to the National Academy of Sciences (USA).

\textbf{2018:} - Shaw Prize in Mathematical Sciences (shared with
Alexander Beilinson) --- \emph{``for their contributions to algebraic
geometry, representation theory, and mathematical physics, in particular
the geometric Langlands program.''}

\textbf{2023:} - Wolf Prize in Mathematics --- \emph{``for contributions
to algebraic geometry, number theory, and mathematical physics,
including the theory of quantum groups and the geometric Langlands
program.''}

\subsection{Honorary Degrees and
Memberships}\label{ch1990vladimirdrinfeld:honorary-degrees-and-memberships}

\begin{itemize}
\tightlist
\item
 Honorary Doctorate, University of Chicago.
\item
 Foreign Member, Royal Society (elected 2024).
\item
 Member, American Academy of Arts and Sciences.
\end{itemize}

\begin{center}\rule{0.5\linewidth}{0.5pt}\end{center}

\section{Selected Major
Publications}\label{ch1990vladimirdrinfeld:selected-major-publications}

\begin{enumerate}
\def\labelenumi{\arabic{enumi}.}
\item
 \textbf{``Elliptic modules.''} \emph{Matematicheskii Sbornik} 94,
 594--627, 1974. --- \emph{The introduction of Drinfeld modules, which
 revolutionized the arithmetic of function fields.}
\item
 \textbf{``Langlands conjecture for \(GL(2)\) over function fields.''}
 \emph{Functional Analysis and its Applications} 10, 1976. ---
 \emph{The first complete proof of a non-abelian case of the Langlands
 conjectures.}
\item
 \textbf{``Hopf algebras and the Yang--Baxter equation.''}
 \emph{Doklady Akademii Nauk SSSR} 283, 1060--1064, 1985. --- \emph{The
 Drinfeld double construction and fundamental results on quantum
 groups.}
\item
 \textbf{``Quantum groups.''} \emph{Proceedings of the International
 Congress of Mathematicians}, Berkeley, 798--820, 1986. --- \emph{The
 well-known ICM address that introduced quantum groups to the
 mathematical world.}
\item
 \textbf{``Lie algebras and equations of KdV type.''} (with V. Sokolov)
 \emph{Journal of Soviet Mathematics} 30, 1975--2036, 1984. ---
 \emph{The Drinfeld--Sokolov reduction and its applications to
 integrable systems.}
\item
 \textbf{``A new realization of Yangians and quantized affine
 algebras.''} \emph{Doklady Akademii Nauk SSSR} 296, 13--17, 1987. ---
 \emph{The representation theory of Yangians and their structure.}
\item
 \textbf{``Quantum groups and the Yang--Baxter equation.''}
 \emph{Zapiski Nauchnykh Seminarov LOMI} 155, 18--49, 1986. ---
 \emph{Further developments in the theory of quantum groups.}
\item
 \textbf{``Chtoukas and the Langlands program.''} \emph{Functional
 Analysis and its Applications} 21, 107--122, 1987. --- \emph{The
 introduction of chtoukas and their role in the Langlands program.}
\item
 \textbf{``Geometric Langlands correspondence.''} (with A. Beilinson)
 \emph{University of Chicago Preprint}, 1991. --- \emph{The
 foundational paper on the geometric Langlands program.}
\item
 \textbf{``On the structure of quasi-triangular quasi-Hopf algebras.''}
 \emph{Journal of Mathematical Sciences} 47, 2631--2656, 1989. ---
 \emph{The theory of quasi-Hopf algebras and their classification.}
\item
 \textbf{``Finite-dimensional representations of Yangians.''}
 \emph{Journal of Mathematical Sciences} 47, 2600--2631, 1989. ---
 \emph{The classification of finite-dimensional Yangian
 representations.}
\end{enumerate}

\begin{center}\rule{0.5\linewidth}{0.5pt}\end{center}

\section{References and Further
Reading}\label{ch1990vladimirdrinfeld:references-and-further-reading}

\subsection{Official Sources}\label{ch1990vladimirdrinfeld:official-sources}

\begin{itemize}
\tightlist
\item
 IMU Fields Medal 1990 Announcement:
 https://www.mathunion.org/imu-awards/fields-medal/fields-medal-1990
\item
 Shaw Prize Announcement, 2018.
\item
 Wolf Prize Citation, 2023.
\end{itemize}

\subsection{Key Research Papers}\label{ch1990vladimirdrinfeld:key-research-papers}

\begin{itemize}
\tightlist
\item
 Drinfeld, V. ``Quantum groups.'' \emph{Proceedings of the ICM} 1986,
 Berkeley, 798--820.
\item
 Drinfeld, V. ``Elliptic modules.'' \emph{Matematicheskii Sbornik}
 94(4), 594--627, 1974.
\item
 Drinfeld, V. ``Langlands conjecture for \(GL(2)\) over function
 fields.'' \emph{Functional Analysis and its Applications} 10(3),
 183--194, 1976.
\item
 Drinfeld, V. ``Hopf algebras and the Yang--Baxter equation.''
 \emph{Doklady Akademii Nauk SSSR} 283(5), 1060--1064, 1985.
\end{itemize}

\subsection{Textbooks and Monographs}\label{ch1990vladimirdrinfeld:textbooks-and-monographs}

\begin{itemize}
\tightlist
\item
 Chari, V. and Pressley, A. \emph{A Guide to Quantum Groups}. Cambridge
 University Press, 1994.
\item
 Frenkel, E. \emph{Langlands Correspondence for Loop Groups}. Cambridge
 University Press, 2007.
\item
 Goss, D. \emph{Basic Structures of Function Field Arithmetic}.
 Springer, 1996.
\item
 Joseph, A. \emph{Quantum Groups and Their Primitive Ideals}. Springer,
 1995.
\item
 Drinfeld, V. ``Langlands' conjecture and the theory of Drinfeld
 modules.'' \emph{Proceedings of the ICM}, Kyoto, 1990.
\end{itemize}

\subsection{Biography and Exposition}\label{ch1990vladimirdrinfeld:biography-and-exposition}

\begin{itemize}
\tightlist
\item
 Manin, Y. ``The Work of Vladimir Drinfeld.'' ICM 1990 Laudatio, Kyoto.
\item
 Beilinson, A. ``Drinfeld's Work on the Geometric Langlands Program.''
 \emph{Current Developments in Mathematics}, 1995.
\item
 Frenkel, E. ``The Work of Drinfeld and the Geometric Langlands
 Program.'' \emph{Notices of the American Mathematical Society}, 41(9),
 1994.
\item
 Bernstein, J. and Zelevinsky, A. ``The Mathematics of Vladimir
 Drinfeld.'' \emph{Shaw Prize Citation}, 2018.
\end{itemize}

\begin{center}\rule{0.5\linewidth}{0.5pt}\end{center}

Drinfeld's work continues to shape large areas of mathematics and
theoretical physics. The Langlands program, quantum groups, and
noncommutative geometry all bear his deep imprint, and his ideas remain
at the forefront of research across multiple disciplines.

\emph{(End of report)}
